\section{Representativeness}
\label{sec:representativeness}

The generators category requires representativeness to be quantified. We
evaluate the building channel against independent real meter corpora, the
behavioral channel against the calibration corpora under a held-out protocol
with bootstrap confidence intervals, and the DER channels against an external
reference implementation. All evaluations use the clean noise profile and raw
(non-peak-normalized) simulated magnitudes.

\subsection{Building load versus real meters}
\label{sec:loadfidelity}

We compare generated load against two independent references: NREL ComStock
end-use load profiles (climate zones 5B, 3B, 4A, 6A) and nineteen real meters
from the Building Data Genome~2 corpus, matched by archetype and size.
Following ASHRAE Guideline~14 we report $\cvrmse{}$ and $\nmbe{}$, together
with normalized load-shape correlation, peak-hour error, and load factor.
Generated \emph{shapes} track the real stock well: shape correlation ranges
$0.71$--$0.94$ across archetypes and peak-hour error is at most three hours.
Absolute intensity reflects the single-prototype scope: four of five
archetypes sit 37--54\% below the stock-average ComStock intensity---as
expected when one code-compliant 2019 prototype represents a heterogeneous
building stock---while the large-office archetype sits 27\% above
($\nmbe{} = -27.1\%$). We report these per archetype (Table~\ref{tab:g14})
rather than in aggregate, and return to the implication in
\S\ref{sec:limitations}.

\begin{table}[t]
\caption{Building-load fidelity versus ComStock and BDG2 (ASHRAE G14 metrics),
per archetype. \owed{Final table from validate\_buildingload via WS-F.}}
\label{tab:g14}
\centering\small
\framebox[\linewidth]{\rule{0pt}{2.2cm} Table 3 placeholder}
\end{table}

\subsection{Behavioral fidelity: held-out and interval-quantified}
\label{sec:behfidelity}

For each calibrated population we generate sessions across seeds and compare,
region by region and variable by variable, against the source corpus: mean
absolute error of the mean ($|\Delta\mu|$), the two-sample KS statistic, and
the Wasserstein-1 distance, with the arrival--dwell copula compared by
Spearman correlation gap. With thousands of sessions per cell the KS
$p$-value is uninformative; we therefore judge by effect size and quantify
uncertainty by a seeded bootstrap over source sessions
(\owed{95\% CIs, WS-A}). Aggregate results: mean $|\Delta\mu|$ of $0.31$\,h
across all region$\times$variable cells; worst-cell KS $0.233$; maximum
copula $\rho$-gap $0.226$. Under the held-out protocol---distributions fitted
on a training split and evaluated against unseen sessions---the median
degradation is $\ks_{\text{holdout}} - \ks_{\text{train}} = 0.012$,
indicating no systematic overfit, though individual dwell cells degrade by up
to $+0.29$ and are disclosed in the full table
(\owed{Tab.~2 with CIs and the F3 protocol note, WS-A/WS-F}).

An ablation isolates the calibration decisions: replacing the pooled arrival
fit broadcast to all regions with per-region fits reduced the largest ACN
region's arrival KS from $0.179$ to $0.079$; admitting mixture families
(\owed{aggregate arrival KS improvement, regenerated via WS-F}) captures the
bimodality of Fig.~\ref{fig:arrival}.

\subsection{DER channels}
\label{sec:derfidelity}

The PV model is deterministic physics, so we validate it against an external
reference implementation---NREL's PVWatts~v8 via the SAM SDK---under
identical configuration and the \emph{same} weather file, eliminating
weather-source confounds. Annual energy agrees to $+1.27\%$ (ours
164.5\,MWh vs.\ 162.4\,MWh for a 100\,kW$_{\mathrm{DC}}$ system); every
month agrees within $\pm 2.6\%$; hourly agreement is well inside ASHRAE
Guideline~14 calibration tolerances ($\cvrmse{} = 15.1\%$,
$\nmbe{} = +1.3\%$, Pearson $r = 0.994$), and a lag scan confirms exact
time-axis alignment. A derate-equalized secondary run isolates the pure
physics difference (isotropic vs.\ Perez transposition, NOCT thermal model)
at $-2.8\%$; both readings sit comfortably under the 5\% acceptance gate.
DR event \emph{timing} follows program rules by construction (season,
notification lead, caps); event \emph{magnitudes}
\owed{are grounded in published program data, or carried as an explicitly
stylized prior --- WS-D outcome}, and we state the resolved position in
\S\ref{sec:limitations}.
